\section{Verification, provenance, and proof boundaries}
\paragraph{Completed mathematical claims.}
The nearest-square lemma and the paired/nested witness transformations
are proved above over all unrestricted integer assignments. The new
Lucas slice lemma treats negative parameters and the exceptional
indices explicitly. Their composition with the pinned source theorem
gives completed consistency equations. They are not merely bounded
search results. The affine-flow theorem, disjunctive scratch reuse,
and asymmetric parity theorem are also unbounded mathematical claims,
but their new global consistency constructors are not finished.

\paragraph{Finite checks are evidence of different kinds.}
The archive checker first verifies its sealed SHA-256 file manifest and
retained-member provenance. In a fresh temporary copy it executes only
the newly authored builders and checkers. This prevents reproduction
from modifying the delivered package. Exact generated outputs must
rebuild byte for byte. Restricted AST readers are used for inherited
polynomial strings; incoming Python is never imported. The independent
source34 checker reconstructs the actual four compiler coefficients.
The direct audit re-expands every local sum of squares and checks
triangular, nonnegative constructor restoration. The displayed compact
low-degree residuals have their own exact symbolic equality check.
Finite signed, Pell-congruence, flow and modular-circuit examples are
regressions, not substitutes for the unbounded arguments.

\paragraph{Negative controls.}
The release includes a manifest tamper test and mathematical mutation
controls in the independent component audits. These have different
meanings: rejecting changed bytes establishes packaging integrity,
whereas rejecting a changed coefficient or invalid flow assumption
exercises a mathematical check. Neither should be described as a
formal kernel certificate. Optimized-Python runs are included for
independent auditors using explicit exceptions, to make clear that
those validations do not disappear under \texttt{python -O}.

\paragraph{Formal status reviewed at the pin.}
The latest repository work includes assembled source-to-ratio-to-finalizer
bridges, a generic-ring evaluator, 31-to-16 elimination, and arithmetic
source-equivalence work for an actual 89-record source. This progress
must not be conflated with a certificate for a new literal circuit.
The new packet's adapted formula identities and global endpoint still
require the corresponding elaboration and proof-kernel replay. The
pinned formal status documents are retained as inert evidence.

\paragraph{Why the fixed integers are not a loophole.}
The coefficient tuple is produced from a finite, explicitly specified
proof-search source, not selected according to its unknown halting
behavior. Payload hashes bind exact integers to exact source conventions.
A large coefficient is counted as a fixed leaf in the operation ledger,
not pretended to be free of specification cost. Raw and compressed sizes
are separately reported. Different compiler-port conventions, especially
shifted and native masks, are never silently interchanged.

\paragraph{Next decisive experiment.}
The most promising departure from the present finalizer is to elaborate
\texttt{SharedAffineProductEquation.lean}, prove its local source bridge,
then materialize and count the global polynomial produced by the direct
bounded-forall compiler. That would test whether replacing ordered
execution by commutative proof balance helps after compilation, rather
than only at the local level. The common multiple, sequence CRT scale,
and internal compiler majorants must remain separate until an explicit
noncircular sharing theorem is proved. A local 17- or 23-witness count
alone says nothing about the eventual nine- or ten-variable frontier.

A second experiment is integrating the asymmetric parity theorem with
a low-cost binary-base guard. The missing binary-base implication is a
specific obstacle, not a generic appeal to DPRM. Neither experiment is
reported as a completed smaller equation here.
